\chapter{Application of Knowledge}
\label{ch:application}

\emph{What does the application of knowledge part of the practical ask, and how should a candidate prepare for it?}
The practical has a part called application of knowledge. The Board names it and gives it marks, but its papers say nothing more about it. This short chapter states exactly what is known, what candidates report, and how to prepare: by learning to choose, adapt and explain practices for a given person.

\begin{unverified}
This chapter is written from a notes block whose fact check is still pending. The marks and the place in the practical follow the Board's guideline; everything else is what candidates report or a study suggestion. Confirm the Board's own format with the Board or an accredited institute.
\end{unverified}

\section{What the Board says}

Application of knowledge carries 10 marks within the practical, which has 140 marks in all. The Board's syllabus and guideline name this part and say the practical includes it; they do not describe it further. The practical assessment of each candidate lasts at most one hour, at the assessor's discretion within that, and this part is assessed inside it.

\section{What candidates report}

Candidates report questions that ask them to apply what they know: choosing and giving a short sequence for a given group, naming the cautions for a condition, and explaining meditation techniques and the stages of yoga nidra. The groups they name include students, women, older people, people with diabetes, high blood pressure, high cholesterol or back pain, and people with a disability. These are reports, not official questions.

\section{How to answer}

Use what you know about the protocol, the practices and the learner to decide what to teach, and give the reason. For any person:
\begin{enumerate}
\item Give the practices that suit the person, in the protocol order.
\item Name the practices to avoid, with the reason.
\item Keep to prevention and well-being; make no claim of cure.
\item Say when to refer to a doctor.
\end{enumerate}
The guidebook describes the skills behind such answers as those of a teacher: applying the protocol, tailoring the instruction to the person, adapting the lesson and communicating clearly. The Board does not say that this part tests them, but they are the natural preparation.

\ycbfig{ycb1-fig-application-of-knowledge}

\section{Preparing}

Work out a few real cases and practise saying a five-minute plan for each aloud: what you would keep, change or leave out, and why. Useful cases are a beginner, an older adult, a person with a stiff back, a person with high blood pressure and a pregnant woman. The practical chapters give the cautions for each practice, and \cref{ch:lessons} gives the method of adapting a plan.

\subsection*{Worked examples}
\begin{itemize}
\item \textbf{A person with high blood pressure.} Keep: loosening practices done slowly, gentle standing and sitting asanas, shavasana, abdominal breathing, nadi shodhana and bhramari without holding the breath, meditation. Leave out: kapalabhati, the inverted posture, strong back bends, and breath holding. Practise only with the doctor's agreement.
\item \textbf{A person with a stiff lower back.} Keep: gentle loosening of the trunk, supported forward bends with bent knees, sarala bhujangasana, makarasana and shavasana. Leave out: deep forward bends with straight legs, vakrasana and strong twists if painful. Refer to a doctor if the pain is severe or persistent.
\item \textbf{Model answer: a short sequence for school students.} Prayer; lively loosening practices; tadasana, vrikshasana, trikonasana; bhujangasana and makarasana; shavasana; bhramari; a short silent sitting; the peace chant. Keep it short and playful.
\end{itemize}

\subsection*{Check yourself}
\begin{enumerate}
\item How many marks does application of knowledge carry, out of how many for the practical?
\item What do the Board's papers say about its content?
\item Give the four points of a good answer for any person.
\item Name three groups that candidates report being asked about.
\end{enumerate}

\subsection*{Answers}
\begin{enumerate}
\item Ten, out of 140 for the practical.
\item Only its name and marks, and that it is part of the practical.
\item Practices that suit, in order; practices to avoid, with reasons; prevention, not cure; when to refer to a doctor.
\item Any three of: students, women, older people, people with diabetes, high blood pressure, high cholesterol, back pain, a disability.
\end{enumerate}

\paragraph{Sources.} Board Guideline for Yoga Professionals, versions 0.2 and 0.3 (marks, duration, components of the practical); the Board syllabus for Level 1; the official guidebook for Level I (2016, older scheme), teaching competences; candidate study notes for what candidates report (not an authority).
